\documentclass[11pt,a4paper]{article}
\usepackage[margin=25mm]{geometry}
\usepackage{amsmath,amssymb,amsthm,hyperref}
\newtheorem{theorem}{Theorem}
\newtheorem{lemma}[theorem]{Lemma}
\newtheorem{proposition}[theorem]{Proposition}
\newcommand{\gen}[1]{\langle#1\rangle}
\title{Further leaves and elementary halving criteria}
\author{Human Compression project: N014 continuation}
\date{}
\begin{document}\maketitle
We use the complete parity-cover theorem and its notation. Thus $\Gamma$
is a residually connected geometry with a simple vertex--edge graph $Q$
on types $0,1$, satisfying (B2): an edge is incident with a higher object
exactly when both its endpoints are. For a higher flag $F$, $Q_F$ is
its common vertex--edge shadow. It is connected and contains an edge.
Write $D(Q)$ for the bipartite double cover, and choose a component $C$.
The halved geometry has the vertices of $C$ as its two new types;
a higher object is $(x,C_x)$, where $C_x$ is a component of $D(Q_x)$
inside $C$. Incidence is membership for a lifted vertex and nonempty
component intersection for incident higher objects. The parity-cover
theorem proves flag extension without an additional parity hypothesis.
It proves residual connectedness precisely under parity compatibility.
Those full proofs are retained inputs, not inferred from this note's
finite example.

Two further elementary consequences recover the graph and duality
clauses of the source. If $Q$ has girth $g$, the girth of its double
cover is the smaller of the least even cycle length in $Q$ and twice
the least odd cycle length, omitting a term when that kind of cycle
does not exist. Indeed an even cycle lifts to a cycle of its own
length, and an odd cycle traversed twice lifts to a cycle of twice
its length. Conversely a shortest cover cycle projects to a closed
walk without immediate reversals. If the projection is simple, it
is an even cycle. If it is not simple, cutting at repeated vertices
either exhibits an even cycle no longer than the walk, or exhibits
at least two odd cycles, counted with multiplicity. The connecting
paths contribute only additional length. This gives the claimed
lower bound. Thus, if $g$ is even, the new gonality is $g/2$; if
$g$ is odd, it is $\min\{g,k/2\}$, where $k$ is the least even
cycle length when one exists. This is source Lemma~3.5.

When $Q$ is nonbipartite, $C=D(Q)$ and the involution
$(v,\epsilon)\mapsto(v,1-\epsilon)$ exchanges the new types.
It sends a higher component to its image under the same involution,
preserving all incidences. It is the duality of Proposition~3.15.

\section{The second leaf}
Suppose that two further types $k,l$ also satisfy (B1) and (B2), with
type $k$ as vertices. Write the two endpoints of a type-$l$ object $y$
as $x_1,x_2$. Conditions (B2) at this leaf give
\[
                 Q_y=Q_{x_1}\cap Q_{x_2}
\]
as vertex--edge subgraphs, even though $\{x_1,x_2\}$ is not a flag.
In particular $Q_y$ is contained in each endpoint shadow.

\begin{theorem}\label{leaf}
The halved geometry satisfies (B1) and (B2) at $(k,l)$ if and only if
every $y$ with bipartite $Q_y$ has at least one endpoint $x_j$ with
bipartite $Q_{x_j}$.
\end{theorem}
\begin{proof}
Fix a lift $(y,C_y)$. Since $C_y$ is connected and is a subgraph of
each $D(Q_{x_j})$, it lies in a unique component $C_{x_j}$. These are
exactly its two lifted endpoints. An endpoint component meeting $C_y$
must be that component; hence there are always exactly two endpoints.
Two distinct lifted edges with the same endpoints must project to the
same $y$, by original (B1). Such a pair exists exactly when $Q_y$ is
bipartite, giving two components, while both $Q_{x_j}$ are nonbipartite,
so both components are contained in the same endpoint lifts. If an
endpoint shadow is bipartite, its two components separate the two lifts
of $y$. This proves that the displayed condition is equivalent to (B1).

Assume the condition, and consider (B2). Any object incident with
$(y,C_y)$ is incident with both its endpoints: for a lifted vertex this
follows from $C_y\subseteq C_{x_j}$; for a higher object it follows
from the same inclusion and original (B2).

Conversely, let an object be incident with both endpoints. First let
it be a lifted vertex $v^\epsilon$. Original (B2) puts $v$ in $Q_y$.
If $Q_y$ is nonbipartite, its only component contains both lifts of
every vertex, so $v^\epsilon\in C_y$. Otherwise choose an endpoint
with bipartite shadow. The restriction of its selected component to
$D(Q_y)$ is exactly $C_y$, so the same conclusion follows.

Now let the object be $(z,C_z)$ of a further type. Original (B2)
gives $z*y$, so $Q_{y,z}$ is connected and contains an edge. If $Q_y$
is nonbipartite, $C_y=D(Q_y)$; since $C_z$ projects onto every vertex
of $Q_z$, it meets $C_y$ over a vertex of $Q_{y,z}$. Otherwise choose
a bipartite endpoint shadow $Q_x$. The selected components $C_z,C_x$
meet by the given incidence. Their restrictions to the connected
graph $Q_{z,x}$ are compatible: $C_x$ selects one bipartite-cover
component, and $C_z$ either selects that component or contains both.
Consequently their intersection projects onto every vertex of
$Q_{z,x}$. Restrict to $Q_{y,z}\subseteq Q_{z,x}$. The component
selected by $C_x$ there is the restriction of $C_y$, so $C_z$ meets
$C_y$. This proves the converse incidence, and hence (B2).
\end{proof}

If an automorphism preserving the types exchanges $x_1,x_2$ while
fixing $y$, their shadows have the same bipartiteness. In particular
this holds in a flag-transitive original geometry. In that case the
criterion reduces to the one printed in Proposition~3.17 of
Piedade--Tranchida: along every incidence $x*y$, either both shadows
are bipartite or neither is. In general it suffices that one endpoint
is bipartite. The necessity argument in the printed proof implicitly
identifies the bipartiteness of the two endpoint shadows.

\subsection*{A finite example distinguishing the two criteria}
Let $V=\mathbb Z/5\mathbb Z$, and let $Q$ be its cycle, with edges
$e_i=\{i,i+1\}$. Use four types. Types $0,1$ are $V,E(Q)$.
Type $2$ consists of $X$ and $x_j$ for $j\in V$; type $3$ consists
of $y_j$ for $j\in V$. The vertex shadow of $X$ is $V$, and those
of $x_j,y_j$ are $V\setminus\{j\}$. A vertex or edge is incident
with one of these objects exactly when it is contained in its shadow.
The only incidences between types $2,3$ are
\[
                          X*y_j,\qquad x_j*y_j.
\]
The two truncations $(0,1)$ and $(2,3)$ are respectively the cycle
and a star with five edges. They satisfy (B1). They also satisfy (B2):
the induced subgraph on $V\setminus\{j\}$ is the path $Q-j$, and
the intersection of the shadows of $X,x_j$ is that same path.

Every flag extends to a chamber. If a type-$3$ object, or a type-$2$
object other than $X$, is specified, the remaining vertex and edge
can be chosen in the nonempty path $Q-j$; a specified vertex has
a neighbor in that path. Complete the higher types with $y_j$ and
either $X$ or $x_j$, respecting any prescribed object. If neither
is specified, choose any compatible cycle
edge and then any $j$ outside its two endpoints. These choices also
cover flags already containing an edge.

Here is a complete check of residual connectedness. For a flag of
types $0,1$, the remaining incidence graph is the star at $X$ with
its incident edges subdivided by the available $y_j$. For a flag
$(v,X)$, the two incident cycle edges are joined through any $y_j$
with $j\notin\{v-1,v,v+1\}$; every other available $y_j$ is
incident with at least one of the two edges. For $(v,x_j)$ the
remaining graph is a star at $y_j$. Flags of types $0,3$ give
$K_{d,2}$ with $d=1$ or $2$; flags of types $1,2$ give $K_{2,3}$
or $K_{2,1}$; flags of types $1,3$ give $K_{2,2}$. Finally a flag
of types $2,3$ has the vertex--edge incidence graph of $Q-j$ as
its residue. These are all six possible pairs of types.

In a residue at a vertex $v$, the object $X$ joins all remaining
edges, and every remaining $x_j,y_j$ is incident with one of them.
In a residue at an edge its two endpoints and $X$ form a connected
core to which every remaining object attaches. Residues at $X,x_j$
or $y_j$ have the connected vertex--edge core $Q$ or $Q-j$, and
the other objects attach to it. The whole geometry has the core
$Q$. Thus all residues of rank at least two are connected.

Now $Q_X=C_5$ is nonbipartite, whereas every $Q_{x_j}=Q_{y_j}$
is the bipartite path on four vertices. The printed condition fails
at each $X*y_j$, but Theorem~\ref{leaf} applies using endpoint $x_j$.
The halved geometry therefore still satisfies both conditions at
$(2,3)$. It is also residually connected: every higher flag other
than $\{X\}$ has a bipartite constituent, so the parity condition
is automatic. The object counts before and after halving are
$(5,5,6,5)$ and $(5,5,11,10)$, respectively. Both geometries have
$60$ chambers. A separate literal incidence program checks all
$322$ original and $362$ transformed flags. The preceding proof
does not rely on that program. The example is not thin or regular;
it addresses precisely the general assertion of Proposition~3.17
and the corresponding clause of Theorem~3.18.

\section{Stabilizers and the induced neighbor graph}
Let $\Gamma$ now be a regular leaf hypertope, with distinguished
involutions $a,b,r_i$, where $a$ commutes with every $r_i$. Fix a
base chamber, its vertex $p$, its edge $e$, and its other endpoint
$q=ap$. Put
\[
 P=G_p=\gen{b,r_i},\qquad R=G_{p,e}=\gen{r_i}.
\]
All groups and actions below are the actual faithful geometric actions.

\begin{proposition}\label{stabilizer}
Condition (B1) holds exactly when $P\cap aPa=R$. If it fails,
the marked halving subgroup $\gen{aba,b,r_i}$ fails the intersection
property.
\end{proposition}
\begin{proof}
The pointwise stabilizer of $p,q$ is $P\cap aPa$. It acts
transitively on the edges joining them. Indeed flag transitivity
sends $(p,e)$ to $(p,e')$; its element fixes $p$ and sends the
other endpoint $q$ to itself. The stabilizer of the base edge in
this action is $R$. Thus the set of edges with these endpoints
is in bijection with $(P\cap aPa)/R$. Every edge has two distinct
endpoints by thinness and the leaf condition, so (B1) is equivalent
to this index being one. All pairs of adjacent vertices are
equivalent under the group, so the base pair suffices.

For the new marking its two endpoint omission subgroups are
$H_0=P$ and $H_1=aPa$, while $H_{0,1}=R$. A strict intersection
therefore violates the marked intersection property.
\end{proof}

For a base object $x_i$ of a higher type put $K_i=G_{x_i}$ and
$P_i=P\cap K_i$. Identify vertices with left cosets $gP$, so
the neighbors of $p=P$ form $PaP/P$. This convention is fixed
throughout the next proposition.

\begin{proposition}\label{B2}
Assume (B1). Condition (B2) is equivalent to
\[
                   K_iP\cap PaP=P_i aP\qquad(i\ne0,1).
\]
The literal sufficient condition in source Lemma~4.6,
\[
 Pa\cap wP=\varnothing\quad
 (w\in K_i,\ wP\ne aP),
\]
implies (B2), but is in fact impossible for a regular leaf hypertope
of rank at least three satisfying (B1).
\end{proposition}
\begin{proof}
The vertices incident with $x_i$ are $K_iP/P$. Since $a\in K_i$,
the edges at $p$ incident with $x_i$ have their other endpoints
$P_i aP/P$: the group $P_i$ is transitive on the appropriate
vertex--edge--object flags. Original (B1) identifies an edge
uniquely with its endpoints. The displayed equality therefore
says exactly that every global neighbor of $p$ in the vertex
shadow of $x_i$ is a neighbor within its residue. This is (B2)
at $p,x_i$; flag transitivity gives it at every vertex--object
incidence. The forward implication of (B2) also follows directly
from the leaf digons.

The printed condition forces any coset in the left-hand set to
be $aP$: an intersection with $PaP$ can be translated on the
right within $P$ to give an intersection with $Pa$. Hence it is
a sufficient, stronger condition. To see why it cannot hold here,
take $w=ba\in K_i$, which is allowed for any higher type $i$.
Then $ba\in Pa\cap baP$. Moreover $baP\ne aP$. Otherwise $b$
would fix both $p,q$, and (B1) would force $b\in R$ by
Proposition~\ref{stabilizer}, contrary to the original intersection
property. Thus this $w$ contradicts the printed condition. The
coset equality above gives the nonvacuous criterion.
\end{proof}

\section{The two elementary examples}
The regular polytope of type $\{3,3,2\}$ has group $S_4\times C_2$,
with the three adjacent transpositions on four letters and an independent
central involution as its generators. Its vertices, edges and triangles
are those of a tetrahedron, and it has two distinct facets, each incident
with all those faces. The intersection of the two facets has several
maximal triangles and no greatest proper face, so its face poset is
not a lattice. Nevertheless the vertex--edge graph is $K_4$, and
incidence with a triangle or either facet is exactly containment of
the endpoints. Hence (B1) and (B2) hold. Higher flags are chains,
whose common vertex shadow is the smallest shadow, so parity
compatibility holds and the repaired halving theorem applies.
This proves the assertion of source Example~4.8 without assuming
that its face poset is a lattice.

For Example~4.9, realize the hemicube by identifying antipodal
faces of the ordinary cube. It has four vertices and six edges,
with graph $K_4$, so (B1) holds. Each of its three faces is a
four-cycle using all four vertices but only four of the six edges.
An omitted edge has both endpoints on that face, proving failure
of (B2). Dually, each pair of its three faces has two common edges,
so the opposite leaf fails (B1).

The automorphism group is the signed permutation group in dimension
three modulo its central sign change, of order $24$, acting faithfully
as $S_4$ on the four antipodal vertex pairs. Choose the usual markers
$a=z_1,b=s_{12},c=s_{23}$: $a$ acts as a double transposition and
$b,c$ as transpositions, with orders $|ab|=4$, $|bc|=3$ and
$[a,c]=1$. More explicitly, choose the representatives
\[
 A=(1,1,1),\ B=(1,-1,-1),\ C=(-1,1,-1),\ D=(-1,-1,1)
\]
of the four antipodal pairs. Then $a=(AB)(CD)$, $b=(BC)$,
$c=(CD)$ and $aba=(AD)$. The new markers $aba,c,b$ are transpositions. Their
successive products have order three and the two endpoints commute.
The subgroup $\gen{b,c}$ is $S_3$ fixing one vertex, while $aba$
joins that vertex to another. They are therefore the adjacent
transpositions of a path on four vertices and generate $S_4$ with
its simplex marking. Their coset geometry is the tetrahedron.
This is an algebraic halving statement; it does not apply the
geometric theorem after its (B2) hypothesis has failed.

\section{Attribution and limits}
The definitions and targets are those of Piedade--Tranchida,
\emph{Constructing new geometries: a generalized approach to halving
for hypertopes}, arXiv:2405.19050v1, especially Propositions~3.17,
4.7, Lemmas~4.5--4.6 and Examples~4.8--4.9. The stabilizer result
and examples are their conclusions. The endpoint criterion and
the explicit coset equality provide the corrections and reusable
arguments supplied here. The fixed preprint's general second-leaf
necessity assertion is refuted; its regular flag-transitive
specialization remains valid. These are source corrections, not
solutions of named open conjectures. No statement is made about
an uninspected journal revision, publication priority or Lean
verification. A distinct mathematical review is still required.
\end{document}
